% Preprint draft: MRIxFields 2026 Task 3, team inzva mri. Springer LNCS format (MICCAI
% workshops); arXiv-ready once the remaining sections are written.
%
% Status (2026-10-08): title, abstract and introduction written. Sections 2-8 are a plan: each
% lists its content and where the numbers come from (TODO.md sections, reports/*.pdf). The
% grey \plan blocks are drafting aids; delete them as each section is written.
%
% Numbers: validation phase = challenge leaderboard (reports/architecture_report.tex); test
% phase = Synapse "Final Ranking" (syn72060672, 1 Oct 2026); local analyses =
% reports/error_spectrum_report.tex and TODO.md. Style: no em dashes.
%
% Compile from inside reports/preprint/:  ~/anaconda3/envs/tex/bin/tectonic preprint.tex
\documentclass[runningheads]{llncs}

\usepackage[T1]{fontenc}
\usepackage{amsmath}
\usepackage{amssymb}
\usepackage{graphicx}
\usepackage{booktabs}
\usepackage{xcolor}
\usepackage[colorlinks=true,linkcolor=black,citecolor=black,urlcolor=black]{hyperref}

\definecolor{plan}{HTML}{627D98}
% Drafting aid: what a section will contain, with its sources. Not part of the paper.
\newcommand{\plan}[1]{{\small\color{plan}\noindent\textbf{Plan.}\ #1\par}}

% Alternative titles:
%   Cross-Field Brain MRI Translation: MIM and Conditioning-Based Approaches (the MICCAI talk)
%   A Field-Conditioned U-Net with Masked Pretraining for Any-to-Any Cross-Field Brain MRI
%   Translation
\title{Information, Not Architecture: Conditioning and Masked Pretraining for Any-to-Any
Cross-Field Brain MRI Translation}
\titlerunning{Conditioning and Masked Pretraining for Cross-Field MRI Translation}

\author{Berkin Deniz Kahya\inst{1,2} \and
Racha Badreddine\inst{1,2} \and
Furkan Y\"uceyal\c{c}\i n\inst{2} \and
Ahmet Zelka\inst{2} \and
\.Ilkay \"Oks\"uz\inst{1}}
\authorrunning{B.\,D. Kahya et al.}
% corresponding-author e-mail to add under the institutes
\institute{\.Istanbul Technical University, PIMI Lab, \.Istanbul, T\"urkiye \and
inzva, \.Istanbul, T\"urkiye}

\begin{document}
\sloppy
\maketitle

% ============================================================================ abstract
\begin{abstract}
Translating a brain MRI from one magnetic field strength to another would let low-field and
high-field scanners share downstream tools. Task 3 of the MRIxFields 2026
challenge asks for one model that maps any of five field strengths (0.1, 1.5, 3, 5 and 7\,T) to
any other, for T1-weighted, T2-weighted and T2-FLAIR images, with three paired travelling
volunteers and 1{,}939 unpaired volumes for training. We describe a 2D U-Net conditioned on the
source domain, the target domain and the axial slice position, pretrained by masked image
modelling on the unpaired volumes, given all three contrasts of the source session, and trained
with an SSIM term in the loss; its prediction averages three checkpoints and four flipped
inputs. In the final test phase (20 unseen volunteers) it reached an SSIM of 0.9426, fourth
of 25 teams on the challenge's main metric and sixth on the mean rank over SSIM, nRMSE and LPIPS,
and received a 3rd Place Award. Ablations on the validation leaderboard attribute most of the
gain over a plain U-Net (SSIM 0.8699 to 0.9137) to conditioning ($+0.0278$) and to pretraining on
the unpaired data ($+0.0103$); five alternative backbones, including conditional flow
matching, scored lower. A spectral analysis separates two regimes:
0.1\,T inputs lack most of the fine detail of 7\,T, while T1W and T2W inputs at 3\,T and 5\,T are
already as fine. The final model's output is too smooth at every frequency, yet the missing power is not
coherent with the target, so post-hoc sharpening and a spectral-profile loss both fail to raise
SSIM.
% Code: \url{https://github.com/denizberkin/mrixfields-analysis} (confirm it is public first)

\keywords{Cross-field MRI translation \and Ultra-low-field MRI \and Masked image modelling
\and Conditional U-Net \and Power spectrum}
\end{abstract}

% ======================================================================== introduction
\section{Introduction}

The magnetic field strength of an MRI scanner sets much of what an image can show. Portable
systems at 0.1\,T and below bring imaging to the bedside and to sites without a shielded suite,
at the cost of signal-to-noise ratio and resolution~\cite{sheth2021portable,arnold2023lowfield};
7\,T systems resolve finer anatomy but remain rare~\cite{ladd2018ultrahigh}. A model that renders
a scan as it would appear at another field strength could let tools built for one regime run on
the other, and make data acquired at different fields comparable~\cite{iglesias2023synthsr}.

The MRIxFields 2026 challenge~\cite{mrixfields2026} poses this as image translation over five
field strengths (0.1, 1.5, 3, 5 and 7\,T) and three contrasts (T1W, T2W and T2FLAIR). Its Task~3
is the general case: one model for all 20 ordered field pairs, 60 translations per subject. The
training data are lopsided. Three travelling volunteers were scanned at every field in every
contrast and give the only paired examples (45 co-registered volumes); 1{,}939 further volumes
from 1{,}056 subjects are unpaired, and none of these subjects has two field strengths in one
contrast. Submissions are scored with SSIM~\cite{wang2004ssim}, the main metric, nRMSE and
LPIPS~\cite{zhang2018lpips}, on 17 validation volunteers and, in the final phase, on 20 unseen
volunteers through a Docker submission; the final rank is the mean of the three metric ranks.

Two properties of the data shaped our work. First, the task is two mappings, not one. Radially
averaged power spectra of the volumes follow $P(f)\propto f^{-\alpha}$, as those of natural
images do~\cite{vanderschaaf1996power}. Sources at 0.1\,T fall off much faster than the 7\,T
targets ($\Delta\alpha = \alpha_{\text{source}} - \alpha_{7\,\mathrm{T}}$ between $+1.01$ and
$+1.50$ across contrasts), so the fine detail of a high-field target is largely absent from the
input; T1W and T2W sources at 3\,T and 5\,T are already finer than 7\,T ($\Delta\alpha$ between
$-0.51$ and $-0.06$), so translating among them mostly changes contrast. Transitions that
involve 0.1\,T are 40\% of the task and carry 50\% of our final model's $1-\mathrm{SSIM}$ on
the paired subjects. Second, three paired subjects leave nothing to hold out: all three are
needed for training, so local scores are in-sample and reward fit. On the two models we scored
both ways, local scoring and the validation leaderboard ranked them in opposite orders, and
model selection was done on the leaderboard.

Our model is a 2D U-Net~\cite{ronneberger2015unet} that translates one axial slice and is told
which translation to perform: learned embeddings of the source domain, the target domain and the
slice position are added at the bottleneck and modulate every decoder block through
FiLM~\cite{perez2018film}. The unpaired volumes enter through masked image
modelling~\cite{he2022mae,xie2022simmim} of the whole network before supervised training on the
paired subjects. The input stacks all three contrasts of the source session next to the image
being translated; the loss adds an SSIM term to L1 and LPIPS; and the submitted prediction
averages three consecutive checkpoints~\cite{izmailov2018swa} and four flipped inputs. On the
validation leaderboard these steps raise SSIM from 0.8699 for a plain U-Net to 0.9137. In the
final test phase the model scored an SSIM of 0.9426, fourth of 25 teams, was sixth on the mean
rank, and received a 3rd Place Award.

Most of what else we tried did not help, and we report it with the same scoring. Swin
UNETR~\cite{hatamizadeh2022swinunetr}, a frozen DINOv3 encoder~\cite{simeoni2025dinov3}, a
LeJEPA-pretrained tubelet encoder~\cite{balestriero2025lejepa}, a transformer that splits
attention by frequency band, and a three-stage conditional flow-matching
model~\cite{lipman2023flow} all scored below the conditional U-Net (0.858 to 0.887, against
0.898). Through-plane context, a later checkpoint and a weight soup~\cite{wortsman2022soups}
each moved validation SSIM by at most 0.001. A spectral analysis of the final model accounts for much of the
remaining error: its output is too smooth at every frequency (72\% of the target's power at
0.1--0.2 cycles/pixel), but the missing detail is not coherent with the target, so restoring the
power, by post-hoc sharpening or by a spectral-profile loss, raises the error with it and does
not raise SSIM.

Our contributions are:
\begin{itemize}
  \item a field-conditioned U-Net with masked pretraining on unpaired volumes and multi-contrast
        input, with every step of its development scored on the challenge leaderboard
        (Sects.~\ref{sec:method} and~\ref{sec:results});
  \item a spectral reading of the task that separates a detail-generation regime (from 0.1\,T)
        from a contrast-mapping regime (among 3, 5 and 7\,T) (Sect.~\ref{sec:analysis});
  \item an error-spectrum and coherence analysis of the final model that explains why
        sharpening and spectrum-matching objectives fail on it (Sect.~\ref{sec:spectrum});
  \item negative results for five alternative backbones and formulations and for changes to
        context, schedule, loss and averaging, each scored against its parent
        (Sect.~\ref{sec:other}).
\end{itemize}

% ============================================================== planned sections (to write)
\section{Challenge and Data}\label{sec:data}
\plan{Task 3 definition; 5 fields $\times$ 3 contrasts; training split (3 paired volunteers,
45 volumes; 1{,}939 unpaired volumes, 1{,}056 subjects), validation (17) and test (20)
volunteers; submission format (20 pairs $\times$ 3 contrasts $\times$ subjects, 30-slice slab);
metrics with formulas (nRMSE, SSIM 7$\times$7 uniform window, LPIPS AlexNet) and the
mean-rank ranking. Copy-the-input baseline: 0.836497. Sources: deck slides 3--7,
\texttt{reports/slides\_miccai.tex}. Figure: \texttt{figures/slide\_task\_fields.pdf}.}

\section{Analysis Before Modelling}\label{sec:analysis}
\plan{Spectral slopes per field and contrast ($\Delta\alpha$ table, 30 retrospective volumes per
cell, fit band 0.02--0.2 cycles/pixel; \texttt{spectral/fig\_rapsd\_by\_field\_retro.pdf}).
Error by field strength (0.1\,T: 40\% of transitions, 50\% of $\sum(1-\mathrm{SSIM})$; 7\,T the
easiest target, 3\,T the easiest source; direction makes no difference; shipped model, fixed
harness, Report III Table 7). Residual maps. Measurement: local scoring in-sample; the wrong-plane
harness fixed 2026-09-11 [\S41], so pre-fix local numbers (seed and subject spread, 0.1\,T
weighting) are void and need re-measuring before they can be cited.}

\section{Method}\label{sec:method}
\plan{32-channel 4-level 2D U-Net; joint (contrast, field) domain embeddings for source and
target plus a sinusoidal slice-position token, summed at the bottleneck and applied by FiLM in
every decoder block; zero-initialised residual head [\S9, \S39]. MIM pretraining of the whole
network: 30{,}000 steps over 1{,}939 volumes [\S18, \S21]. 4-channel input: the image being
translated plus T1W, T2W, T2FLAIR at the source field [\S23, \S25]. Loss
$L_1 + 0.5\,(1-\mathrm{SSIM}) + 0.05\,\mathrm{LPIPS}$, SSIM matched to the scorer [\S28].
Inference: mean of epochs 6--8 weights, 4 flips [\S40]; background zeroed by the source mask.
Implementation details: optimiser, schedule, batch, hardware, training time (from the run
configs and audit logs).}

\section{Results}\label{sec:results}
\plan{Validation ladder (table: vanilla 0.869856, + conditioning 0.897646, + FiLM \& residual
0.892839, + MIM 0.903110, + 4 channels 0.906773, + SSIM loss 0.908873, + axial token 0.909433,
+ averaging \& TTA 0.913652, with nRMSE and LPIPS). Final test ranking, top 9 ordered by SSIM:
0.942551 (4th), nRMSE 0.248030 (6th), LPIPS 0.050566 (14th); mean rank 8.00, 6th tied. The
SSIM/LPIPS trade-off of the last three steps (SSIM $+0.0069$, LPIPS $+0.0101$;
\texttt{figures/showcase/trade\_*.pdf}). Qualitative 3$\times$5 grids up from 0.1\,T and down
from 7\,T with residuals (\texttt{figures/showcase/grid\_*}).}

\section{Error Spectrum of the Final Model}\label{sec:spectrum}
\plan{Error power over target power by band and source field; cumulative error (47\% below 0.05
cycles/pixel); direction does not separate the spectra; error on folds (35\% of squared error on
18\% of pixels). Prediction/target power below 1 in every band; SSIM factor split (structure
largest). Coherence identity $C = (P_{\hat y} + P_y - P_e)/2$. Post-hoc sharpening (17 settings,
no gain) and the $\Delta\alpha$-gated spectral-profile loss (power 0.743 $\to$ 0.882, coherence
0.833 $\to$ 0.822, SSIM $-0.00118$, 164/168 transitions lower). Source:
\texttt{reports/error\_spectrum\_report.pdf}, [\S44].}

\section{Other Routes}\label{sec:other}
\plan{Backbones and formulations with validation SSIM: Swin UNETR 3D 0.858214, LeJEPA tubelet +
UNETR 0.859000, DINOv3 ViT 0.878918, FPS-Former 0.879531, conditional flow matching 0.886725
(\texttt{figures/showcase/other\_backbones.pdf}). Other changes against their parent run,
validation SSIM or fixed-harness local (Report III Table 17, \texttt{other\_changes\_4.pdf}).
Synthetic pairs from degraded 7\,T: never trained on; the reason they were dropped is a void
pre-fix result. Workshop methods of the other top teams for context.}

\section{Discussion and Conclusion}\label{sec:discussion}
\plan{What the results say: information (which pair, unpaired data, other contrasts, the metric
in the loss) bought the gains; architecture did not. Limitations: three paired subjects, all
qualitative figures in-sample, local scoring coarser than the leaderboard, LPIPS traded for
SSIM. Future work: fastMRI brain (6{,}970 volumes, 1.5 and 3\,T, raw multi-coil
k-space~\cite{zbontar2018fastmri}): spectral slopes at scale, coherence per frequency for
reconstruction, noise floors before model comparison.}

\bibliographystyle{splncs04}
\bibliography{references}

\end{document}
